\begin{figure}[htbp]
    \raggedright
    \resizebox{0.99\textwidth}{!}{%
    \begin{tikzpicture}[
        font=\large,
        node distance=5mm,
        box/.style={draw, rounded corners=2mm, minimum width=3.35cm,
                    minimum height=1.8cm, text width=3.15cm,
                    align=center, thick, fill=blue!7},
        output/.style={draw, rounded corners=2mm, minimum width=3.35cm,
                    minimum height=1.8cm, text width=3.15cm,
                    align=center, thick, fill=green!10},
        arrow/.style={-{Latex[length=3mm]}, thick}
    ]
        \node[box] (acq) {Adquisición\\EEG, MEG, fMRI o ECoG};
        \node[box, right=of acq] (prep) {Preproceso\\filtrado, artefactos y normalización};
        \node[box, right=of prep] (repr) {Representación\\temporal, espacial o tiempo--frecuencia};
        \node[box, right=of repr] (dec) {Modelo de decodificación\\clasificación o generación};
        \node[output, right=of dec] (out) {Salida\\texto, voz o control de dispositivo};

        \draw[arrow] (acq) -- (prep);
        \draw[arrow] (prep) -- (repr);
        \draw[arrow] (repr) -- (dec);
        \draw[arrow] (dec) -- (out);
        \coordinate (feedbackleft) at ($(acq.north)+(0,9mm)$);
        \coordinate (feedbackright) at (out.north |- feedbackleft);
        \draw[arrow, rounded corners=2mm] (out.north) -- (feedbackright)
            -- node[above, align=center, font=\normalsize]
            {realimentación\\en sistemas en línea}
            (feedbackleft) -- (acq.north);
    \end{tikzpicture}}
    \caption{Etapas generales de una interfaz cerebro--computador. La salida del decodificador puede utilizarse directamente o cerrar el bucle mediante realimentación al usuario.}
    \label{fig:bci_pipeline}
\end{figure}
